\title{色彩空间与显示设备色彩再现}
\author{杨子凡}
\date{Jun 20, 2026}
\maketitle
当同一张照片在手机、显示器和打印件上呈现出明显不同的色彩时，我们其实遭遇的是设备间「色彩方言」的冲突。每一台显示设备都拥有自己独特的色彩再现能力，这种能力由其物理结构与驱动算法共同决定，而人眼感知到的颜色则依赖于光谱功率分布、物体表面反射率以及视觉匹配函数的共同作用。CIE 1931 XYZ 色彩空间正是将这种主观感知转化为客观坐标的数学框架，它为后续的色彩管理提供了公共参照。\par
\chapter{回到原点：光源、物体、观察者}
可见光谱涵盖了约 380 nm 至 780 nm 的电磁波段，而三原色理论指出，绝大多数颜色可由红、绿、蓝三种单色光按不同比例混合得到。颜色在数学上可表示为光谱功率分布与物体反射率以及人眼匹配函数的积分结果，这一过程在 CIE 1931 标准观察者模型中被形式化。CIE 1931 XYZ 由此将人眼响应转换为三个刺激值，使得任何可见色都可以在三维坐标系中被唯一标识。\par
\chapter{色彩空间的「地图」——从 CIE 到设备}
CIE 1931 与 1976 色度图描绘了人眼所能感知的全部颜色范围，而实际显示设备只能覆盖其中的有限子集。sRGB、Adobe RGB、DCI-P3 及 Rec.2020 等色彩空间正是通过定义各自的原色坐标与白点，将 CIE 全地图「裁剪」成适合自身硬件的封闭区域。色域（Gamut）在 CIE 色度图上表现为一个封闭多边形，其面积大小直接反映设备能够再现的色彩范围。进一步引入亮度维度后，色域体积（Gamut Volume）则描述了设备在三维空间中可呈现的颜色总量。\par
\chapter{显示设备如何「画」颜色？}
显示设备通过背光或自发光光源产生初始光谱，再经彩色滤光片与子像素布局调制后形成最终颜色。LED、Mini-LED、量子点、OLED 与激光光源各自具有不同的光谱功率分布，直接影响设备原色的色坐标。原色坐标决定了设备色域的三个顶点位置，而亮度分布与白场稳定性则制约了高动态范围内容的重现能力。HDR 峰值亮度与白场漂移现象，正是这些物理限制在实际使用中的具体表现。\par
\chapter{再现过程的「三重门」——Tone、Gamut、EOTF}
当高动态范围场景映射到有限亮度显示器时，Tone Mapping 负责将场景亮度压缩至设备可接受范围。Gamut Mapping 则处理源色彩空间与目标色彩空间之间的范围差异，通过压缩或扩展算法完成颜色转换。电光转换函数（EOTF）及其逆函数定义了数字编码值与显示亮度之间的非线性关系，PQ、HLG 与 sRGB gamma 曲线分别针对不同应用场景优化。3D LUT 以离散查找表方式实现任意映射，而矩阵加曲线方法则依赖线性变换与一维曲线组合，二者在计算复杂度与精度上各有权衡。\par
在 sRGB 色彩空间中，gamma 曲线可表示为线性段与幂函数的拼接，公式为\par
$$ L = \begin{cases} \frac{V}{12.92} & V \leq 0.04045 \\ \left( \frac{V + 0.055}{1.055} \right)^{2.4} & V > 0.04045 \end{cases} $$\par
其中 ( V ) 为 0 到 1 归一化的编码值，( L ) 为对应的线性亮度。该公式确保了编码值在低灰阶区域具有更高精度，符合人眼对暗部细节的敏感特性。\par
PQ 曲线则采用感知量化方式，其核心公式为\par
$$ L = \left( \frac{\max\left( \left( \frac{V}{c_1} \right)^{1/m_2} - c_2, 0 \right)}{c_3 - c_4 \left( \frac{V}{c_1} \right)^{1/m_2}} \right)^{1/m_1} $$\par
其中常数 ( c\_{}1, c\_{}2, c\_{}3, c\_{}4, m\_{}1, m\_{}2 ) 由 SMPTE ST 2084 标准定义。该函数能够在 10 bit 编码下覆盖高达 10000 nits 的亮度范围，同时保持人眼可分辨的亮度阶调。\par
\chapter{色彩管理：让「翻译」自动化}
ICC Profile 通过描述设备色彩特性与 PCS（Profile Connection Space）之间的双向转换关系，实现跨设备色彩一致性。色彩管理系统（CMS）的工作流程可概括为输入设备色彩首先转换至 PCS，再由 PCS 转换至输出设备色彩，从而避免直接在设备间进行不可逆的色彩映射。软校色利用显示器模拟印刷结果，硬校色则通过实际打印件验证色彩准确性。DisplayCAL 配合 i1Display Pro 等硬件校色仪，可逐点测量显示器响应并生成包含 gamma、色温与色域信息的「屏幕身份证」，即 ICC Profile 文件。\par
\chapter{真实世界的冲突与解决方案}
HDR 视频在 SDR 显示器上播放时，由于 Tone Mapping 与 Gamut Mapping 同时进行，容易出现色块伪影与细节丢失。RAW 摄影文件经过 sRGB 与 CMYK 多次 gamut 裁剪后，饱和度会显著下降，原因是每次映射都不可逆地丢弃了超出目标色域的颜色。多屏拼接场景中，白点不一致会导致视觉断层，根源在于各显示器在 CIE 色度图上选择了不同的白场坐标。ICtCp 色彩空间、ACEScct 工作流程以及 HDR10+ 与 Dolby Vision 的动态 gamut mapping 算法，正在从标准层面缓解这些实际问题。\par
内容创作者应在整个后期流程中保持统一的工作色彩空间，仅在最终交付时才转换为目标设备空间，并善用视图 Proof 模式检查色彩是否超出目标 gamut。普通用户在启用系统广色域模式前，宜先完成硬件校色，以避免因 gamma 与白点偏差导致的色彩偏移。选购显示器时，应关注其实际覆盖的色域标准而非仅看 NTSC 百分比。未来 Micro-LED 全色发光、超广色域印刷油墨以及 AI 自适应 gamut mapping 技术，将进一步缩小设备间「色彩方言」的鸿沟，让每一块屏幕都能准确说同一种颜色语言。\par
